\section{A convention allowing singleton loops}
\label{c:app-loops}

Suppose a loop is allowed to cover its single incident vertex.
Let \(\ell_i\) be the number of distinguished loop choices at vertex
\(i\), and retain the multiplicities \(k_{\{i,j\}}\) for pairs of
distinct vertices.  All these integers are given in binary.
A matching under this convention covers every vertex exactly once by
either a chosen loop or an ordinary pair.
Write \(N_j\) for the number of such matchings using exactly \(j\)
loops, and \(N_{\mathrm{loop}}=\sum_{j=0}^n N_j\).

\begin{proposition}
\label{c:prop-loops}
There is one deterministic algorithm which, given a graph on \(n\)
vertices with nonnegative binary integer pair and singleton-loop
multiplicities, returns a nonnegative rational number
\(A_{\mathrm{loop}}\) such that
\[
             \frac{N_{\mathrm{loop}}}{2^{18n}}
                  \le A_{\mathrm{loop}}\le N_{\mathrm{loop}}.
\]
It returns zero exactly when \(N_{\mathrm{loop}}=0\).
Its running time and output length are polynomial in the input bit
length, including when the input gives \(n\) in binary and lists only
the positive pair and loop records.
\end{proposition}

\begin{proof}
First count the distinct vertices occurring in a positive ordinary pair
or having a positive loop count. If any of the \(n\) vertices is missing,
it cannot be covered, so return zero. This test uses only the listed
records, without enumerating the vertex set. Otherwise, if there are
\(r\) positive pair records and \(s\) positive loop records, then
\(n\le2r+s\). Thus the constructions indexed by vertices below have
size polynomial in the input length, even for a compact binary vertex
count.

For each \(j\in\{0,\ldots,n\}\), construct an ordinary loopless graph
\(G_j\) by adjoining \(j\) labeled new vertices.
Keep all original edges between distinct vertices.
There are no edges among the new vertices, and the pair joining
original vertex \(i\) to any new vertex has multiplicity \(\ell_i\).
Label these parallel edges by the corresponding original loop choices.

Every perfect matching of \(G_j\) pairs the new vertices to \(j\)
distinct original vertices.  Replace those pairs by their labeled
loop choices and retain the other edges; this gives a matching counted
by \(N_j\).
Conversely, each matching counted by \(N_j\) has exactly \(j!\) lifts,
one for each bijection from the new vertices to its \(j\) loop-covered
vertices.  Edge labels preserve the loop choices independently of this
bijection.  Thus, including \(j=0\),
\begin{equation}
\label{c:eq-loop-count-identity}
                         N(G_j)=j!\,N_j.
\end{equation}

Run the ordinary algorithm on each \(G_j\), obtaining a nonnegative
integer \(A_j\), and
return
\begin{equation}
\label{c:eq-loop-estimate}
                         A_{\mathrm{loop}}
                              =\sum_{j=0}^n\frac{A_j}{j!}.
\end{equation}
Parity is tested on each transformed order \(n+j\), as part of that
ordinary algorithm.  There is no initial rejection for odd \(n\):
a single vertex with a loop is already a positive-count example.
The parity of \(n+j\) equals that of the number \(n-j\) of original
vertices that remain to be paired by ordinary edges.

By Theorem~\ref{thm:algorithm-guarantee} and
Equation~\eqref{c:eq-loop-count-identity},
\[
             \frac{N_j}{512^{\,n+j}}
                   \le\frac{A_j}{j!}\le N_j.
\]
Since \(n+j\le2n\), summing these inequalities proves the claimed
factor \(512^{2n}=2^{18n}\).
All terms are nonnegative, and \(A_j=0\) exactly when \(N_j=0\), so
zero recognition is also exact.

The construction retains compressed multiplicity records:
\(G_j\) has at most \(r+js\) positive pair records. No parallel-edge
list is expanded. Each transformed graph has at most \(2n\) vertices and uses only
multiplicities already present in the input.  If \(b_{\max}\) is the largest
binary exponent of any positive pair or loop multiplicity, taking
\(b_{\max}=0\) when all multiplicities are zero,
Proposition~\ref{d:prop-bit-complexity} bounds each call and gives
\(O(n^2+nb_{\max}+1)\) bits for \(A_j\).
There are only \(n+1\) calls.
The rational sum can be represented exactly with the common denominator
\(n!\):
\begin{equation}
\label{c:eq-loop-common-denominator}
       A_{\mathrm{loop}}
          =\frac{\displaystyle\sum_{j=0}^n
                              A_j\,\frac{n!}{j!}}{n!}.
\end{equation}
The denominator has \(O(n\log_2(n+1)+1)\) bits, and the numerator has
\(O(n^2+nb_{\max}+n\log_2(n+1)+1)\) bits.
Factorials, exact integer divisions, products, and the sum therefore
take polynomial bit time.  The case \(n=0\) uses the single term
\(j=0\), with \(0!=1\) and \(A_0=1\), so the empty matching is counted
exactly.
\end{proof}
